\documentclass{article}
\usepackage{booktabs}
\usepackage{array}
\usepackage{amssymb} % For math symbols
\newcolumntype{L}{>{\raggedright\arraybackslash}p{0.3\textwidth}}
\newcolumntype{R}{>{\raggedright\arraybackslash}p{0.5\textwidth}}
% Define custom commands for clarity and simplicity
\newcommand{\stablecoin}{\textsc{StableCoin}}
\newcommand{\usd}{\textsc{USD}}
\newcommand{\apy}{\textsc{APY}}
\newcommand{\pos}{\textsc{Proof-of-Stake}}
\newcommand{\usdc}{\textsc{USDC}}







\usepackage{amsmath}


% ------------------------------------------------------------------------------
% Required Packages for preamble/font.tex
% ------------------------------------------------------------------------------
\makeatletter
% Font encoding and families
\@ifpackageloaded{fontenc}{}{\usepackage[T1]{fontenc}}   % better font encoding
\@ifpackageloaded{lmodern}{}{\usepackage{lmodern}}       % Latin Modern — scalable replacement for CM

% Typography helpers
\@ifpackageloaded{lettrine}{}{\usepackage{lettrine}}      % dropped capitals
\@ifpackageloaded{lipsum}{}{\usepackage{lipsum}}          % dummy text (testing)
\makeatother

% ------------------------------------------------------------------------------
% Documentation and Usage Examples: lipsum and lettrine
% ------------------------------------------------------------------------------
% DESCRIPTION
%   The lipsum package provides dummy text for testing typographic layout.
%   The lettrine package produces dropped capitals at the beginning of a paragraph.
%
% FEATURES
%   - lipsum: generate one or more paragraphs for quick layout checks
%   - lettrine: control dropped capital lines, spacing, and scaling
%
% USAGE: lipsum
%   Basic:    \lipsum                  % generates a default paragraph
%   Paragraph: \lipsum[1]              % first paragraph
%   Range:     \lipsum[2-3]            % paragraphs 2 to 3
%   Words:     \lipsum[1][1-5]         % words 1–5 from paragraph 1
%
% USAGE: lettrine
%   Basic:     \lettrine{T}{his} starts a paragraph with a dropped “T”
%   Lines:     \lettrine[lines=3]{T}{his}    % drop across 3 lines
%   Spacing:   \lettrine[findent=2pt,nindent=0pt]{T}{his}
%   Scaling:   \lettrine[lines=3,loversize=0.15]{T}{his}
%   Slope:     \lettrine[hang=1]{T}{his}     % tighter hang of text under the drop cap
%
% EXAMPLES (copy into your document body to preview)
%   % Example 1: Simple lipsum
%   % \lipsum[1]
%
%   % Example 2: Variable Length lipsum
%   % \lipsum[2-3]
%
%   % Example 3: Drop cap with default settings
%   % \lettrine{L}{orem} ipsum dolor sit amet, consectetur adipiscing elit.\lipsum[2]
%
%   % Example 4: Three-line drop cap with adjusted size
%   % \lettrine[lines=3,loversize=0.2]{D}{olor} sit amet, consectetur adipiscing elit.\lipsum[3]
%
%   % Example 5: Fine-tuned indents for tighter text wrap
%   % \lettrine[lines=3,findent=1.5pt,nindent=0pt]{T}{ypography} testing block.\lipsum[4]
%
% NOTES
%   - Use \par after a lettrine paragraph if spacing looks off in complex layouts.
%   - Pair lettrine with a serif body font for best aesthetics.


\usepackage[margin=1in]{geometry}



\begin{document}
\section*{USDC vs. USDT: A Simple Comparison of Stablecoins}

\noindent Two well-known fiat-backed \stablecoin s in the digital asset market are USDC (USD Coin) and USDT (Tether). Both are made to keep their value equal to the U.S. dollar ($\$1.00$), which helps provide stability in the often unpredictable cryptocurrency world.

\subsection*{Definition and Purpose}

\noindent Stablecoins, like USDC and USDT, are used to reduce price 
swings\footnote{Price swings, also called volatility, mean that the 
value of something can change quickly and unpredictably.}. They make 
it possible to send money quickly at any time and act as a temporary 
way to hold value (a \inlinequote{digital dollar}) 
for traders and investors, without needing regular banks.

\subsection*{Key Differences in Simple Terms}

\noindent The main differences between USDC and USDT are in how their 
reserves\footnote{Reserves are the assets held to back up the value of 
the stablecoin, making sure it can be exchanged for real dollars.}, 
openness\footnote{Openness, also called transparency, means how much 
information is shared about how the stablecoin works and what backs it.}, 
and rules from governments\footnote{Regulatory engagement means how 
closely the stablecoin follows laws and rules set by governments.}. 
These differences affect how much people trust each stablecoin, as 
shown in Table \ref{tab:comparison}.

\begin{table}[h!]
    \centering
    \setlength{\tabcolsep}{6pt}
    \renewcommand{\arraystretch}{1.5}
    \begin{tabular}{@{} >{\bfseries}L R @{}}
        \toprule
        \textbf{Feature} & \textbf{Simple Explanation} \\
        \midrule
        \textbf{Transparency} & USDC shares regular reports from independent 
        companies about its reserves. USDT has been questioned in the past 
        for not sharing as much information. \\
        \textbf{Reserve Backing} & USDC is backed mostly by \usd\ and safe 
        U.S. government bonds\footnote{U.S. Treasury instruments are loans 
        to the U.S. government, considered very safe.}. USDT uses a wider 
        mix of assets, including cash, loans, and other investments. \\
        \textbf{Regulatory Standing} & USDC tries to follow U.S. laws 
        closely, and its company (Circle) is based in the U.S. USDT has 
        faced some issues in the past about whether its reserves are 
        enough and if it follows all rules. \\
        \textbf{Market Adoption} & USDT is used the most worldwide and 
        has the biggest market size because it has been around longer. 
        USDC is chosen more by people and companies who want a stablecoin 
        that follows more rules. \\
        \bottomrule
    \end{tabular}
    \caption{Comparison of USDC and USDT Characteristics}
    \label{tab:comparison}
\end{table}

\subsection*{Conclusion}

\noindent Choosing between these two \stablecoin s depends on what matters most to you:

\begin{itemize}
    \item If you want \textbf{clear information and strong rules}, USDC is usually seen as more trustworthy.
    \item If you need \textbf{the most trading options and popularity}, USDT is the most widely used.
\end{itemize}



\newpage

\section*{What does ``USDC \(\cdot\) 4.10\% APY'' mean?}

When a platform (e.g., Coinbase) displays
\[
\text{USDC}\;\cdot\;4.10\%\ \text{APY}
\]
it means: if you hold USDC (a U.S.-dollar-pegged stablecoin) on that platform, you may earn interest at an \textbf{Annual Percentage Yield (APY)} of \(\,4.10\%\) (as of the displayed rate). APY is the effective annual return after accounting for compounding.

\bigskip

\subsection*{Key definitions}
\begin{itemize}
  \item \textbf{USDC:} a stablecoin roughly pegged 1:1 to USD (so 1 USDC \(\approx\) \$1).
  \item \textbf{APY (Annual Percentage Yield):} the effective annual rate including the effect of compounding.
  \item \textbf{APR (Annual Percentage Rate):} usually quoted without compounding; APY converts APR to an annualized, compounding-aware rate.
\end{itemize}

\subsection*{Mathematics of compounding}
Let:
\[
P = \text{initial principal (USD)}, \qquad r = \text{nominal annual rate (decimal)}, \qquad n = \text{number of compounding periods per year}.
\]
If interest is compounded \(n\) times per year, the balance after one year is
\[
A = P\left(1+\frac{r}{n}\right)^{n}.
\]
The corresponding APY is
\[
\text{APY} = \left(1+\frac{r}{n}\right)^{n} - 1.
\]
If interest were compounded continuously, the limit gives
\[
A_{\text{cont}} = Pe^{r},\qquad \text{APY}_{\text{cont}} = e^{r}-1.
\]

\bigskip

\subsection*{Numeric example with \(\mathbf{4.10\%}\) APY}
Take \(P=\$1{,}000\) and the displayed APY \(=4.10\%\), i.e. \(r_{\text{APY}}=0.041\).

\paragraph{Simple annual compounding (once per year).}
If the APY comes from annual compounding (\(n=1\)), then
\[
A = 1000\cdot(1+0.041)=1000\cdot 1.041 = \$1041.00.
\]

\paragraph{Daily compounding approximation (\(n=365\)).}
If interest is compounded daily and the nominal/quoted rate leads to the same effective APY,
\[
A = 1000\left(1+\frac{0.041}{365}\right)^{365}.
\]
Numerically,
\[
A \approx 1000\cdot\left(1+\frac{0.041}{365}\right)^{365} \approx \$1041.85.
\]

\paragraph{Continuous compounding (for comparison).}
\[
A_{\text{cont}} = 1000\,e^{0.041}\approx \$1041.85.
\]

(You can see daily and continuous compounding give virtually identical results at this rate; the small differences are due to the limit behavior.)

\bigskip

\subsection*{Practical notes and caveats}
\begin{itemize}
  \item Platforms may display an APY that is \emph{variable} and can change over time depending on market conditions and the platform's yield sources.
  \item The actual way the platform generates yield (e.g., lending, market-making, staking, partner programs) and any fees or restrictions are important to check in the platform's terms.
  \item Cryptocurrency holdings are exposed to platform risk (counterparty risk) and regulatory risk; APY is not the same as an FDIC-insured bank interest rate.
  \item Many platforms compound daily (or credit earnings daily), but check the platform documentation for the exact compounding frequency and whether there are any minimum balance or withdrawal limits.
\end{itemize}

\bigskip

\section*{Short summary (one line)}
\begin{align*}
    &\text{USDC}\;\cdot\;4.10\%\ \text{APY} \\
    &\qquad\Rightarrow\quad \text{an effective annual return of } 4.10\% \\
    &\qquad\qquad\text{i.e. } \$1000 \to \approx \$1041.85\ \text{after 1 year (with daily/continuous compounding).}
\end{align*}



\newpage

\section*{Crypto Staking, Fees, and Secure Local Storage: A Quick Guide}


\section*{Explanation of the Staking Reward ($\text{Up to } 15.12\% \text{ APY}$)}

\noindent The notice \inlinequote{Earn up to 15.12\% APY $\cdot$ Buy and 
stake eligible assets to earn rewards} promotes \textbf{Staking}, 
which is a high-yield method of earning passive income on 
cryptocurrency, distinct from the rewards earned on simply holding \usdc.

\subsection*{What is Staking?}

\noindent Staking is the process by which a cryptocurrency holder commits their tokens to support the operations of a blockchain network utilizing the \pos\ consensus mechanism.

\begin{itemize}
  \item \textbf{Function:} The staked assets are used to validate transactions and secure the network, replacing energy-intensive mining.
  \item \textbf{Reward Source:} Rewards are paid by the blockchain network itself, directly to the stakers, in the form of the asset being staked. Coinbase handles the technical requirements of this process.
\end{itemize}

\subsection*{The Annual Percentage Yield (APY)}

\noindent The rate of $\text{up to } 15.12\% \text{ APY}$ represents the potential 
annualized return on the highest-yielding eligible asset available for staking on the platform (e.g., Cosmos (\texttt{ATOM})).

\begin{itemize}
  \item \textbf{Nature of Rate:} The rate is \textbf{variable} and fluctuates based on the amount of crypto being staked network-wide and the network's activity. It is not a guaranteed fixed return.
  \item \textbf{Asset-Specific:} Each cryptocurrency eligible for staking will have its own unique \apy\ (e.g., Ethereum (\texttt{ETH}) and Solana (\texttt{SOL}) offer lower rates than $15.12\%$).
\end{itemize}

\subsection*{Key Difference from USDC Rewards}

\noindent Staking offers a higher potential reward but introduces additional financial risk and liquidity constraints compared to holding \usdc\ for rewards.

\begin{table}[h!]
  \centering
  \setlength{\tabcolsep}{6pt}
  \renewcommand{\arraystretch}{1.5}
  \caption{Comparison of Staking vs. USDC Rewards}
  \begin{tabular}{@{} >{\bfseries}L L L @{}}
    \toprule
    \textbf{Category} & \textbf{Staking (e.g., ATOM)} & \textbf{USDC Rewards} \\
    \midrule
    \textbf{APY} & High ($\text{up to } 15.12\%$), \textbf{Variable} & Lower ($4.10\%$), \textbf{Stable} \\
    \textbf{Asset Risk} & \textbf{High.} Value of the staked coin can drop drastically. & \textbf{Low.} Value is pegged to $\$1.00$. \\
    \textbf{Liquidity} & \textbf{Restricted.} Assets are locked; an unstaking period (hours to weeks) is required. & \textbf{Liquid.} Assets are immediately available for trading/sending once the ACH hold clears. \\
    \bottomrule
  \end{tabular}
\end{table}


\newpage




\section*{1. How Staking Works}

\subsection*{Definition}
\textbf{Staking} is the process of locking up a certain amount of cryptocurrency in order to support the operations of a blockchain network that uses a \textbf{Proof-of-Stake (PoS)} or related consensus mechanism. In return, stakers earn rewards (similar to interest) for helping validate transactions and maintain network security.

\subsection*{Mechanics}
In PoS-based blockchains (e.g., Ethereum, Cardano, Solana):
\begin{itemize}
  \item Users lock tokens in a staking contract.
  \item Validators (or staking pools) are chosen to propose and verify blocks in proportion to their stake.
  \item Successful validation yields a reward distributed to participants.
\end{itemize}

If a validator acts maliciously or stays offline, part of their stake may be \emph{slashed} (lost as a penalty).

\subsection*{Mathematical View}
Let:
\[
P = \text{amount staked}, \quad r = \text{annual staking yield (decimal)}, \quad t = \text{time in years}.
\]
Then the staking reward (ignoring compounding) is approximately
\[
R = P \cdot r \cdot t.
\]
If rewards are compounded continuously,
\[
A = P e^{r t}.
\]

\subsection*{Notes}
\begin{itemize}
  \item Staking is available only for PoS-based assets (e.g., ETH, ADA, SOL), not for Bitcoin or XRP.
  \item Rewards depend on network conditions, validator performance, and lock-up duration.
  \item On platforms like Coinbase, staking is custodial --- the platform stakes on behalf of users and takes a small commission.
\end{itemize}

\bigskip

\section*{2. Crypto Transaction Fees: Network vs. Platform (Coinbase)}

\subsection*{(a) Network Fees}
Network fees are payments made directly to blockchain validators or miners for processing and confirming transactions. These fees depend on:
\begin{itemize}
  \item \textbf{Blockchain congestion:} When the network is busy, fees rise.
  \item \textbf{Transaction complexity:} Larger or multi-step transactions require higher computation and storage.
  \item \textbf{Token type:} Each blockchain (Bitcoin, Ethereum, Solana, etc.) has its own fee structure.
\end{itemize}

Mathematically, a simplified transaction fee can be written as
\[
\text{Fee}_{\text{network}} = \text{GasUsed} \times \text{GasPrice}.
\]
For example, on Ethereum, \texttt{GasUsed} measures computational effort, and \texttt{GasPrice} (in gwei) represents the cost per unit.

\subsection*{(b) Platform Fees (Coinbase)}
When trading or transferring on Coinbase, two types of additional fees may apply:
\begin{itemize}
  \item \textbf{Trading spread:} A small percentage (typically $0.5\%$) between buy and sell prices.
  \item \textbf{Coinbase fee:} A service charge depending on payment method, transaction size, or region.
\end{itemize}
These are separate from blockchain-level network fees.  
Coinbase automatically deducts the appropriate network fee when you withdraw crypto to an external wallet.

\bigskip

\section*{3. Keeping Coins Locally and Securely}

\subsection*{Custodial vs. Non-Custodial Storage}
\begin{itemize}
  \item \textbf{Custodial:} The platform (e.g., Coinbase) holds your private keys. You access coins through your account but do not control the keys.
  \item \textbf{Non-custodial:} You hold your own private keys through a wallet (hardware or software). This provides full control but also full responsibility.
\end{itemize}

\subsection*{Local Storage Options}
\begin{enumerate}
  \item \textbf{Hardware Wallets:} Physical devices (e.g., Ledger, Trezor) that store private keys offline. Considered the most secure for long-term holding.
  \item \textbf{Software Wallets:} Apps (e.g., Coinbase Wallet, MetaMask) that store private keys locally on your phone or computer.
  \item \textbf{Paper Wallets:} Physical printouts of keys or QR codes. Secure from hacking but can be lost or damaged.
\end{enumerate}

\subsection*{Security Best Practices}
\begin{itemize}
  \item Always enable \textbf{2FA (Two-Factor Authentication)} on exchange accounts.
  \item Never share or upload your \textbf{seed phrase} (12 or 24-word recovery phrase).
  \item Keep offline backups of your seed phrase in multiple secure locations.
  \item Be cautious of phishing sites and emails pretending to be from Coinbase.
  \item Consider using a dedicated hardware wallet for large holdings.
\end{itemize}

\subsection*{Summary Table}

\begin{center}
\begin{tabular}{|l|l|l|}
\hline
\textbf{Storage Type} & \textbf{Who Holds Keys?} & \textbf{Security Level} \\
\hline
Exchange (e.g., Coinbase) & Platform & Medium \\
Software Wallet (e.g., MetaMask) & You & High (if secured) \\
Hardware Wallet (e.g., Ledger) & You & Very High \\
\hline
\end{tabular}
\end{center}

\bigskip

\section*{Key Takeaways}
\begin{itemize}
  \item Staking earns rewards by helping secure PoS networks.
  \item Transaction fees include both network-level and platform-level charges.
  \item Local storage gives you full control, but also full responsibility for security.
\end{itemize}




\end{document}